\documentclass[10pt,a4paper]{amsart}
\usepackage{amsmath,amssymb,amsthm,graphicx}
% Retained minimal preamble for the Stacks Project chapter varieties.tex.
% Verbatim excerpt of the chapter preamble: Xy-pic setup (native preamble lines 14-18), the native theorem
% declarations with the mathematical macros (native preamble lines 154-217), and the project own no-op
% definitions of the reference/slogan/history display environments (native wrapper.tex lines 1-3) that the
% counted statement uses for its slogan. Copied mechanically; nothing added, removed or redefined.
% For commutative diagrams we use Xy-pic
\usepackage[all]{xy}

% We use 2cell for 2-commutative diagrams.
\xyoption{2cell}
% Theorem environments.
%
\theoremstyle{plain}
\newtheorem{theorem}[subsection]{Theorem}
\newtheorem{proposition}[subsection]{Proposition}
\newtheorem{lemma}[subsection]{Lemma}

\theoremstyle{definition}
\newtheorem{definition}[subsection]{Definition}
\newtheorem{example}[subsection]{Example}
\newtheorem{exercise}[subsection]{Exercise}
\newtheorem{situation}[subsection]{Situation}

\theoremstyle{remark}
\newtheorem{remark}[subsection]{Remark}
\newtheorem{remarks}[subsection]{Remarks}

\numberwithin{equation}{subsection}

% Macros
%
\def\lim{\mathop{\mathrm{lim}}\nolimits}
\def\colim{\mathop{\mathrm{colim}}\nolimits}
\def\Spec{\mathop{\mathrm{Spec}}}
\def\Hom{\mathop{\mathrm{Hom}}\nolimits}
\def\Ext{\mathop{\mathrm{Ext}}\nolimits}
\def\SheafHom{\mathop{\mathcal{H}\!\mathit{om}}\nolimits}
\def\SheafExt{\mathop{\mathcal{E}\!\mathit{xt}}\nolimits}
\def\Sch{\mathit{Sch}}
\def\Mor{\mathop{\mathrm{Mor}}\nolimits}
\def\Ob{\mathop{\mathrm{Ob}}\nolimits}
\def\Sh{\mathop{\mathit{Sh}}\nolimits}
\def\NL{\mathop{N\!L}\nolimits}
\def\CH{\mathop{\mathrm{CH}}\nolimits}
\def\proetale{{pro\text{-}\acute{e}tale}}
\def\etale{{\acute{e}tale}}
\def\QCoh{\mathit{QCoh}}
\def\Ker{\mathop{\mathrm{Ker}}}
\def\Im{\mathop{\mathrm{Im}}}
\def\Coker{\mathop{\mathrm{Coker}}}
\def\Coim{\mathop{\mathrm{Coim}}}

% Boxtimes
%
\DeclareMathSymbol{\boxtimes}{\mathbin}{AMSa}{"02}

%
% Macros for moduli stacks/spaces
%
\def\QCohstack{\mathcal{QC}\!\mathit{oh}}
\def\Cohstack{\mathcal{C}\!\mathit{oh}}
\def\Spacesstack{\mathcal{S}\!\mathit{paces}}
\def\Quotfunctor{\mathrm{Quot}}
\def\Hilbfunctor{\mathrm{Hilb}}
\def\Curvesstack{\mathcal{C}\!\mathit{urves}}
\def\Polarizedstack{\mathcal{P}\!\mathit{olarized}}
\def\Complexesstack{\mathcal{C}\!\mathit{omplexes}}
% \Pic is the operator that assigns to X its picard group, usage \Pic(X)
% \Picardstack_{X/B} denotes the Picard stack of X over B
% \Picardfunctor_{X/B} denotes the Picard functor of X over B
\def\Pic{\mathop{\mathrm{Pic}}\nolimits}
\def\Picardstack{\mathcal{P}\!\mathit{ic}}
\def\Picardfunctor{\mathrm{Pic}}
\def\Deformationcategory{\mathcal{D}\!\mathit{ef}}
\newenvironment{reference}{}{}
\newenvironment{slogan}{}{}
\newenvironment{history}{}{}
\usepackage{cleveref}
\crefname{originallemma}{original lemma}{original lemmas}
\Crefname{originallemma}{Original lemma}{Original lemmas}
\makeatletter
\newlabel{lemma-m-regular-extend-base-field}{{08A4}{source}{}{sourcepointer-1}{}}
\newlabel{lemma-m-regular-extend-base-field@cref}{{[originallemma][0][]08A4}{[1][0][]source}{}{}{}}
\makeatother
\crefname{originalreference}{original reference}{original references}
\Crefname{originalreference}{Original reference}{Original references}
\makeatletter
\newlabel{coherent-lemma-flat-base-change-cohomology}{{02KH}{source}{}{sourcepointer-2}{}}
\newlabel{coherent-lemma-flat-base-change-cohomology@cref}{{[originalreference][0][]02KH}{[1][0][]source}{}{}{}}
\makeatother
\crefname{originallemma}{original lemma}{original lemmas}
\Crefname{originallemma}{Original lemma}{Original lemmas}
\makeatletter
\newlabel{lemma-exact-sequence-induction}{{08A0}{source}{}{sourcepointer-3}{}}
\newlabel{lemma-exact-sequence-induction@cref}{{[originallemma][0][]08A0}{[1][0][]source}{}{}{}}
\makeatother
\crefname{originalremark}{original remark}{original remarks}
\Crefname{originalremark}{Original remark}{Original remarks}
\makeatletter
\newlabel{remark-exact-sequence-induction-cohomology}{{0EGK}{source}{}{sourcepointer-4}{}}
\newlabel{remark-exact-sequence-induction-cohomology@cref}{{[originalremark][0][]0EGK}{[1][0][]source}{}{}{}}
\makeatother
\crefname{originallemma}{original lemma}{original lemmas}
\Crefname{originallemma}{Original lemma}{Original lemmas}
\makeatletter
\newlabel{lemma-hyperplane}{{089Z}{source}{}{sourcepointer-5}{}}
\newlabel{lemma-hyperplane@cref}{{[originallemma][0][]089Z}{[1][0][]source}{}{}{}}
\makeatother
\title{For an m-regular coherent sheaf the multiplication map from the m-th twist and linear forms onto the (m+1)-st twist is surjective}
\author{The Stacks Project Authors}
\date{stacks-project:varieties.tex:local-snapshot-20261011; GNU FDL 1.2 or later}
\begin{document}
\maketitle
\noindent\textit{Editorial note.} All delivered mathematical text is the author's own, copied verbatim from the pinned Stacks Project chapter file varieties.tex; no statement or proof step was rewritten, reordered or removed.\par\medskip
\noindent\textit{Editorial note.} The counted span cites five cross-chapter labels (including a remark), each recorded in assembly.external\_references with its canonical tag as a navigation pointer.\par\medskip
\noindent\textit{Editorial note.} The proof cites no bibliography item inside the counted span.\par\medskip
\noindent\textit{Allowed standard prerequisites (editorial).}\begin{itemize}
\item The induction step producing a hyperplane restriction of a coherent sheaf. Stacks Project, Varieties, Lemma lemma-exact-sequence-induction (tag 0ABH) and Remark remark-exact-sequence-induction-cohomology (tag 0ABI). Supplies the short exact sequence on which the induction step is based.
\item Base-field invariance of m-regularity and flat base change for cohomology. Stacks Project, Varieties, Lemma lemma-m-regular-extend-base-field (tag 0A2A); Cohomology of Schemes, Lemma coherent-lemma-flat-base-change-cohomology (tag 02KH). Justify the reduction to an infinite field and the identification of the base-changed multiplication map.
\end{itemize}\medskip
\ifdefined\hypertarget\hypertarget{sourcepointer-1}{}\fi
\noindent\textit{Source reference.} Original lemma 08A4: sources/stacks-chow/tags.txt, source lines 10519--10519. Cross-chapter standard prerequisite cited by the author inside the counted span; it is used as a pointer and its own proof is not part of this item.\par\medskip
\ifdefined\hypertarget\hypertarget{sourcepointer-2}{}\fi
\noindent\textit{Source reference.} Original reference 02KH: sources/stacks-chow/tags.txt, source lines 3372--3372. Cross-chapter standard prerequisite cited by the author inside the counted span; it is used as a pointer and its own proof is not part of this item.\par\medskip
\ifdefined\hypertarget\hypertarget{sourcepointer-3}{}\fi
\noindent\textit{Source reference.} Original lemma 08A0: sources/stacks-chow/tags.txt, source lines 10515--10515. Cross-chapter standard prerequisite cited by the author inside the counted span; it is used as a pointer and its own proof is not part of this item.\par\medskip
\ifdefined\hypertarget\hypertarget{sourcepointer-4}{}\fi
\noindent\textit{Source reference.} Original remark 0EGK: sources/stacks-chow/tags.txt, source lines 18100--18100. Cross-chapter standard prerequisite cited by the author inside the counted span; it is used as a pointer and its own proof is not part of this item.\par\medskip
\ifdefined\hypertarget\hypertarget{sourcepointer-5}{}\fi
\noindent\textit{Source reference.} Original lemma 089Z: sources/stacks-chow/tags.txt, source lines 10514--10514. Cross-chapter standard prerequisite cited by the author inside the counted span; it is used as a pointer and its own proof is not part of this item.\par\medskip
\begin{lemma}
\label{lemma-m-regular-multiply}
Let $k$ be a field. Let $n \geq 0$.
Let $\mathcal{F}$ be a coherent sheaf on $\mathbf{P}^n_k$.
If $\mathcal{F}$ is $m$-regular, then the multiplication map
$$
H^0(\mathbf{P}^n_k, \mathcal{F}(m)) \otimes_k
H^0(\mathbf{P}^n_k, \mathcal{O}(1))
\longrightarrow
H^0(\mathbf{P}^n_k, \mathcal{F}(m + 1))
$$
is surjective.
\end{lemma}

\begin{proof}
Let $k'/k$ be an extension of fields. Let $\mathcal{F}'$
be as in Lemma \ref{lemma-m-regular-extend-base-field}. By
Cohomology of Schemes, Lemma \ref{coherent-lemma-flat-base-change-cohomology}
the base change of the linear map of the lemma to $k'$ is the
same linear map for the sheaf $\mathcal{F}'$. Since $k \to k'$
is faithfully flat it suffices to prove the lemma over $k'$, i.e.,
we may assume $k$ is infinite.

\medskip\noindent
Assume $k$ is infinite. We prove the lemma by induction on $n$.
The case $n = 0$ is trivial as $\mathcal{O}(1) \cong \mathcal{O}$
is generated by $T_0$. For $n > 0$ apply
Lemma \ref{lemma-exact-sequence-induction}
and tensor the sequence by $\mathcal{O}(m + 1)$ to get
$$
0 \to \mathcal{F}(m) \xrightarrow{s} \mathcal{F}(m + 1) \to
i_*\mathcal{G}(m + 1) \to 0
$$
see Remark \ref{remark-exact-sequence-induction-cohomology}.
Let $t \in H^0(\mathbf{P}^n_k, \mathcal{F}(m + 1))$.
By induction the image $\overline{t} \in H^0(H, \mathcal{G}(m + 1))$
is the image of $\sum g_i \otimes \overline{s}_i$ with
$\overline{s}_i \in \Gamma(H, \mathcal{O}(1))$ and
$g_i \in H^0(H, \mathcal{G}(m))$. Since $\mathcal{F}$ is $m$-regular
we have $H^1(\mathbf{P}^n_k, \mathcal{F}(m - 1)) = 0$, hence long
exact cohomology sequence associated to the short exact sequence
$$
0 \to \mathcal{F}(m - 1) \xrightarrow{s} \mathcal{F}(m) \to
i_*\mathcal{G}(m) \to 0
$$
shows we can lift $g_i$ to $f_i \in H^0(\mathbf{P}^n_k, \mathcal{F}(m))$.
We can also lift $\overline{s}_i$ to
$s_i \in H^0(\mathbf{P}^n_k, \mathcal{O}(1))$ (see proof of
Lemma \ref{lemma-hyperplane} for example). After subtracting the
image of $\sum f_i \otimes s_i$ from $t$ we see that we may assume
$\overline{t} = 0$. But this exactly means that $t$ is the
image of $f \otimes s$ for some $f \in H^0(\mathbf{P}^n_k, \mathcal{F}(m))$
as desired.
\end{proof}

\end{document}
